\documentclass[a4paper]{article}
\usepackage{amsfonts}
\usepackage{amsmath}
\usepackage{amssymb}
\usepackage{graphicx}     

\newcommand{\dint}{\displaystyle\int\limits}
\newcommand{\Bgtan}{\operatorname{Bgtan}}

\begin{document}
	
	\noindent \large Auteur: Dina Haajou \\
	\noindent \large Studierichting:  Informatica\\
	\noindent \large Oefeningenreeks 6D nr. 30.14\\
	
	\medskip
	
	\normalsize
	
	$\displaystyle\sum\limits_{n=1}^{\infty} \dfrac{(-1)^n}{\sqrt{n}}$\\
	
	Criterium 13 (Leibniz):\\
	
	$(|x_n|)_n = \left(\dfrac{1}{\sqrt{n}}\right)_n$ is dalend\\
	
	$\displaystyle\lim\limits_{n \to +\infty} \left| \dfrac{(-1)^n}{\sqrt{n}} \right| = 0 \Rightarrow$ convergent\\
	
	$\displaystyle\sum\limits_{n=1}^{\infty} \left| \dfrac{(-1)^n}{\sqrt{n}} \right| = \displaystyle\sum\limits_{n=1}^{\infty} \dfrac{1}{\sqrt{n}}$ divergeert ($p = \dfrac{1}{2} < 1$)\\
	
	$\Rightarrow$ De reeks convergeert, maar de reeks van de absolute waarden divergeert.\\
	
	\begin{thebibliography}{999}
		%\bibitem{a} Referentie
	\end{thebibliography}

\end{document}
